\documentclass[12pt, letterpaper]{article}
\usepackage{graphicx}
\usepackage{amsmath}
\usepackage{amssymb}
\usepackage{float}
\graphicspath{{images/}}
\title{Logistic regression model}
\author{Evert Jan Karman}
\date{January 1900}
\begin{document}
\maketitle

Logistic regression model

\section{Model}

We obtain n data points with d features.
\newline
We convert the output to values -1 and 1. That will be stored in n x 1 matrix \(Y\).
\newline
We may remove the reviews labeled neutral from our dataset (the data with known sentiment)
\newline
In module 2, d features means how often d-th word appears in the original review text
\newline
We insert a dummy column with ones as 0-th column which we'll map to the intercept \(w_0\)
\newline
So we put the features in n x (d+1) matrix \(H\)
\newline
We also have d + 1 coefficients which we put in a (d+1) x 1 matrix \(w\) where \(w_0\) is the intercept
\newline
The scores per data point are obtained by multiplying \(H\) with \(w\):
\newline
\(H \times w\) = score will be a n x 1 matrix
\newline
After that we convert each single score to a probability using sigmoid function:
\[
P = \frac{1}{1 + e^{-score}}
\]


\section{Side note}
So far a data point was a row in an n x (d+1) matrix. But alternatively sometimes a row \(x_i\) is denoted as:
\[
\begin{bmatrix}
h_0(x_i)\\
h_1(x_i)\\
\vdots\\
h_d(x_i)
\end{bmatrix}
\]

Then the score of the i-th data point is subsequently written as \underline{dot product}
\[
score = w^\top h(x_i)
\]

\end{document}
